\ifdefined\gkver\else\def\gkver{student}\fi
\documentclass[\gkver]{gaokaozhenti}
\begin{document}

\chapter{1996年上海}

\begin{enumerate}
\item
%% number: 1
%% paperName: 1996年上海高考第 1 题 4 分
%% typeId: 1
%% type: 单选题
%% chapter: 原子和原子核
%% point: 原子核式结构
%% method: 无
%% score: 4
%% degree: 850
%% duplicateId: 0
%% body:
根据卢瑟福的原子核式结构模型，下列说法中正确的是\xzanswer{D}
\fourchoices[answer=D]
{原子中的正电荷均匀分布在整个原子范围内}
{原子中的质量均匀分布在整个原子范围内}
{原子中的正电荷和质量都均匀分布在整个原子范围内}
{原子中的正电荷和几乎全部质量都集中在很小的区域范围内}
%% answer: D
%% memo:
\memoanswer{
本题原卷及材料中未提供详解，待补充。
}

\item
%% number: 2
%% paperName: 1996年上海高考第 2 题 4 分
%% typeId: 1
%% type: 单选题
%% chapter: 曲线运动
%% point: 平抛运动
%% method: 无
%% score: 4
%% degree: 650
%% duplicateId: 0
%% body:
物体作平抛运动时，描述物体在竖直方向的分速度 $v_{y}$（取向下为正）随时间变化的图线是\xzanswer{D}
\fourchoices[answer=D, ispicture=true, h=3.2cm]
{\includesvg{figs/02a}}
{\includesvg{figs/02b}}
{\includesvg{figs/02c}}
{\includesvg{figs/02d}}
%% answer: D
%% memo:
\memoanswer{
本题原卷及材料中未提供详解，待补充。
}

\item
%% number: 3
%% paperName: 1996年上海高考第 3 题 4 分
%% typeId: 1
%% type: 单选题
%% chapter: 电路
%% point: 电功电功率
%% method: 近似与估算
%% score: 4
%% degree: 850
%% duplicateId: 0
%% body:
一只普通的家用照明白炽灯正常发光时，通过它的电流值与下列数值较为接近的是\xzanswer{C}
\fourchoices[answer=C]
{$20\UA$}
{$2\UA$}
{$0.2\UA$}
{$0.02\UA$}
%% answer: C
%% memo:
\memoanswer{
本题原卷及材料中未提供详解，待补充。
}

\item
%% number: 4
%% paperName: 1996年上海高考第 4 题 4 分
%% typeId: 1
%% type: 单选题
%% chapter: 光学
%% point: 光电效应
%% method: 无
%% score: 4
%% degree: 650
%% duplicateId: 0
%% body:
当某种单色光照射到金属表面时，金属表面有光电子逸出。如果光的强度减弱，频率不变，则\xzanswer{B}
\fourchoices[answer=B]
{光的强度减弱到某一最低数值时，就没有光电子逸出}
{单位时间内逸出的光电子数减少}
{逸出光电子的最大初动能减少}
{单位时间内逸出的光电子数和光电子的最大初动能都要减小}
%% answer: B
%% memo:
\memoanswer{
本题原卷及材料中未提供详解，待补充。
}

\item
%% number: 5
%% paperName: 1996年上海高考第 5 题 4 分
%% typeId: 1
%% type: 单选题
%% chapter: 电磁感应
%% point: 电磁感应电路分析
%% method: 无
%% score: 4
%% degree: 450
%% duplicateId: 0
%% body:
如图所示，两个互相用电阻不计的短导线连接的金属圆环，粗金属环的电阻为细金属环电阻的二分之一，磁场垂直穿过粗金属环所在区域。当磁感强度随时间均匀变化时，在粗环内产生的感应电动势为 $\varepsilon$，则两环连接处 a、b 间的电势差为\xzanswer{C}
\onepicture[w=5.5cm]{\includesvg{figs/05}}
\fourchoices[answer=C]
{$\frac{1}{2}\varepsilon$}
{$\frac{1}{3}\varepsilon$}
{$\frac{2}{3}\varepsilon$}
{$\varepsilon$}
%% answer: C
%% memo:
\memoanswer{
本题原卷及材料中未提供详解，待补充。
}

\item
%% number: 6
%% paperName: 1996年上海高考第 6 题 4 分
%% typeId: 1
%% type: 单选题
%% chapter: 电磁感应
%% point: 楞次定律推广
%% method: 无
%% score: 4
%% degree: 650
%% duplicateId: 0
%% body:
如图所示，MN 是一根固定的通电长直导线，电流方向向上。今将一矩形金属线框 abcd 放在导线上，让线框的位置偏向导线的左边，两者彼此绝缘。当导线中的电流突然增大时，线框整体受力情况为\xzanswer{A}
\onepicture[w=4.5cm]{\includesvg{figs/06}}
\fourchoices[answer=A]
{受力向右}
{受力向左}
{受力向上}
{受力为零}
%% answer: A
%% memo:
\memoanswer{
本题原卷及材料中未提供详解，待补充。
}

\item
%% number: 7
%% paperName: 1996年上海高考第 7 题 4 分
%% typeId: 1
%% type: 单选题
%% chapter: 运动和力
%% point: 牛顿第二定律
%% method: 无
%% score: 4
%% degree: 450
%% duplicateId: 0
%% body:
如图所示，底板光滑的小车上用两个量程为 $20\UN$、完全相同的弹簧秤甲和乙系住一个质量为 $1\Ukg$ 的物块，在水平地面上，当小车作匀速直线运动时，两弹簧秤的示数均为 $10\UN$。当小车作匀加速直线运动时，弹簧秤甲的示数变为 $8\UN$，这时小车运动的加速度大小是\xzanswer{B}
\onepicture[w=7cm]{\includesvg{figs/07}}
\fourchoices[answer=B]
{$2\Umsq$}
{$4\Umsq$}
{$6\Umsq$}
{$8\Umsq$}
%% answer: B
%% memo:
\memoanswer{
当弹簧测力计甲的示数变为 $8\UN$ 时，弹簧测力计乙的示数变为 $12\UN$，这时物块所受的合力为 $4\UN$。由牛顿第二定律 $F = ma$ 得物块的加速度 $a = \frac{F}{m} = 4\Umsq$，故选项 B 正确。
}

\item
%% number: 8
%% paperName: 1996年上海高考第 8 题 4 分
%% typeId: 1
%% type: 单选题
%% chapter: 运动和力
%% point: 牛顿定律的应用
%% method: 无
%% score: 4
%% degree: 450
%% duplicateId: 0
%% body:
某消防队员从一平台上跳下，下落 $2\Um$ 后双脚触地，接着他用双腿弯屈的方法缓冲，使自身重心又下降了 $0.5\Um$。在着地过程中地面对他双脚的平均作用力估计\xzanswer{B}
\fourchoices[answer=B]
{自身所受重力的 2 倍}
{自身所受重力的 5 倍}
{自身所受重力的 8 倍}
{自身所受重力的 10 倍}
%% answer: B
%% memo:
\memoanswer{
本题原卷及材料中未提供详解，待补充。
}

\item
%% number: 9
%% paperName: 1996年上海高考第 9 题 5 分
%% typeId: 2
%% type: 多选题
%% chapter: 气体性质
%% point: 分子动理论
%% method: 无
%% score: 5
%% degree: 650
%% duplicateId: 0
%% body:
下列叙述中正确的是\xzanswer{BC}
\fourchoices[answer=BC]
{物体的内能与物体的温度有关，与物体的体积无关}
{物体的温度越高，物体中分子无规则运动越剧烈}
{物体体积改变，内能可能不变}
{物体在压缩时，分子间存在着斥力，不存在引力}
%% answer: BC
%% memo:
\memoanswer{
本题原卷及材料中未提供详解，待补充。
}

\item
%% number: 10
%% paperName: 1996年上海高考第 10 题 5 分
%% typeId: 2
%% type: 多选题
%% chapter: 电磁感应
%% point: 自感与互感，涡流
%% method: 无
%% score: 5
%% degree: 650
%% duplicateId: 0
%% body:
如图电路（a）、（b）中，电阻 $R$ 和自感线圈 $L$ 的电阻值都很小，接通 K，使电路达到稳定，灯泡 S 发光，则\xzanswer{AD}
\onepicture[w=7cm]{\includesvg{figs/10}}
\fourchoices[answer=AD]
{在电路（a）中，断开 K，S 将渐渐变暗}
{在电路（a）中，断开 K，S 将先变得更亮，然后渐渐变暗}
{在电路（b）中，断开 K，S 将渐渐变暗}
{在电路（b）中，断开 K，S 将先变得更亮，然后渐渐变暗}
%% answer: AD
%% memo:
\memoanswer{
本题原卷及材料中未提供详解，待补充。
}

\item
%% number: 11
%% paperName: 1996年上海高考第 11 题 5 分
%% typeId: 2
%% type: 多选题
%% chapter: 运动和力
%% point: 牛顿第二定律
%% method: 无
%% score: 5
%% degree: 650
%% duplicateId: 0
%% body:
一向右运动的车厢顶上悬挂两单摆 M 与 N，他们只能在图示平面内摆动。某一瞬时出现图示情景，由此可知车厢的运动及两单摆相对车厢运动的可能情况是\xzanswer{AB}
\onepicture[w=5.5cm]{\includesvg{figs/11}}
\fourchoices[answer=AB]
{车厢作匀速直线运动，M 在摆动，N 静止}
{车厢作匀速直线运动，M 在摆动，N 也在摆动}
{车厢作匀速直线运动，M 静止，N 在摆动}
{车厢作匀加速直线运动，M 静止，N 也静止}
%% answer: AB
%% memo:
\memoanswer{
本题原卷及材料中未提供详解，待补充。
}

\item
%% number: 12
%% paperName: 1996年上海高考第 12 题 5 分
%% typeId: 2
%% type: 多选题
%% chapter: 机械波
%% point: 波的多解性
%% method: 无
%% score: 5
%% degree: 450
%% duplicateId: 0
%% body:
一列横波在某时刻的波形图如图中实线所示，经 $2 \times 10^{-2}\Us$ 后的波形如图中虚线所示，则该波的波速 $v$ 和频率 $f$ 可能是\xzanswer{ABD}
\onepicture[w=8cm]{\includesvg{figs/12}}
\fourchoices[answer=ABD]
{$v$ 为 $5\Ums$}
{$v$ 为 $45\Ums$}
{$f$ 为 $50\UHz$}
{$f$ 为 $37.5\UHz$}
%% answer: ABD
%% memo:
\memoanswer{
本题原卷及材料中未提供详解，待补充。
}

\item
%% number: 13
%% paperName: 1996年上海高考第 13 题 5 分
%% typeId: 2
%% type: 多选题
%% chapter: 交流电
%% point: 变压器
%% method: 无
%% score: 5
%% degree: 450
%% duplicateId: 0
%% body:
在绕制变压器时，某人误将两个线圈绕在图示变压器铁芯的左右两个臂上，当通以交流电时，每个线圈产生的磁通量都只有一半通过另一个线圈，另一半通过中间的臂。已知线圈 1、2 的匝数之比 $n_{1}:n_{2} = 2:1$，在不接负载的情况下\xzanswer{BD}
\onepicture[w=5.5cm]{\includesvg{figs/13}}
\fourchoices[answer=BD]
{当线圈 1 输入电压 $220\UV$ 时，线圈 2 输出电压为 $110\UV$}
{当线圈 1 输入电压 $220\UV$ 时，线圈 2 输出电压为 $55\UV$}
{当线圈 2 输入电压 $110\UV$ 时，线圈 1 输出电压为 $220\UV$}
{当线圈 2 输入电压 $110\UV$ 时，线圈 1 输出电压为 $110\UV$}
%% answer: BD
%% memo:
\memoanswer{
本题原卷及材料中未提供详解，待补充。
}

\item
%% number: 14
%% paperName: 1996年上海高考第 14 题 4 分
%% typeId: 3
%% type: 填空题
%% chapter: 电路
%% point: 欧姆定律
%% method: 无
%% score: 4
%% degree: 850
%% duplicateId: 0
%% body:
（A 组）本世纪初，科学家发现，某些金属材料，当\tkanswer[1.6cm]{温度}降低到一个临界值以下时，会出现电阻\tkanswer[3.4cm]{突然变为零}的现象，这种现象叫做超导现象。
%% answer: 
\jdanswer{
温度，突然变为零（或变为零）
}
%% memo:
\memoanswer{
本题原卷及材料中未提供详解，待补充。
}

\item
%% number: 15
%% paperName: 1996年上海高考第 15 题 4 分
%% typeId: 3
%% type: 填空题
%% chapter: 曲线运动
%% point: 人造地球卫星
%% method: 无
%% score: 4
%% degree: 650
%% duplicateId: 0
%% body:
（A 组）已知地球的质量为 $M$，万有引力恒量为 $G$，地球半径为 $R$。用以上各量表示在地球表面附近运行的人造地球卫星的第一宇宙速度 $v =$ \tkanswer[2.6cm]{$\sqrt{\frac{GM}{R}}$}。
%% answer: 
\jdanswer{
$\sqrt{\frac{GM}{R}}$
}
%% memo:
\memoanswer{
本题原卷及材料中未提供详解，待补充。
}

\item
%% number: 16
%% paperName: 1996年上海高考第 16 题 4 分
%% typeId: 3
%% type: 填空题
%% chapter: 磁场
%% point: 磁力矩
%% method: 整体与隔离
%% score: 4
%% degree: 450
%% duplicateId: 0
%% body:
（A 组）如图所示，在光滑水平桌面上，有两根弯成直角的相同金属棒，它们的一端均可绕固定转动轴 O 自由转动，另一端 b 互相接触，组成一个正方形线框，正方形每边长均为 $l$，匀强磁场的方向垂直桌面向下。当线框中通以图示方向的电流 $I$ 时，两金属棒在 b 点的相互作用力为 $f$，则此时磁感强度的大小为\tkanswer[2.4cm]{$\frac{\sqrt{2}f}{Il}$}。（不计电流产生的磁场）
\onepicture[h=5.5cm, align=right]{\includesvg{figs/16}}
%% answer: 
\jdanswer{
$\frac{\sqrt{2}f}{Il}$
}
%% memo:
\memoanswer{
本题原卷及材料中未提供详解，待补充。
}

\item
%% number: 14
%% paperName: 1996年上海高考第 14 题 4 分
%% typeId: 3
%% type: 填空题
%% chapter: 光学
%% point: 全反射
%% method: 无
%% score: 4
%% degree: 850
%% duplicateId: 0
%% body:
（B 组）当光线由光\tkanswer[1.2cm]{密}（填疏或密）媒质射到光\tkanswer[1.2cm]{疏}（填疏或密）媒质的分界面上时，如果入射角大于\tkanswer[1.4cm]{临界}角，就会发生全反射现象。
%% answer: 
\jdanswer{
密、疏、临界
}
%% memo:
\memoanswer{
本题原卷及材料中未提供详解，待补充。
}

\item
%% number: 15
%% paperName: 1996年上海高考第 15 题 4 分
%% typeId: 3
%% type: 填空题
%% chapter: 曲线运动
%% point: 人造地球卫星
%% method: 无
%% score: 4
%% degree: 650
%% duplicateId: 0
%% body:
（B 组）已知地球表面重力加速度为 $g$，地球半径为 $R$，万有引力恒量为 $G$，用以上各量表示，地球质量 $M =$ \tkanswer[2.4cm]{$\frac{R^{2}g}{G}$}。
%% answer: 
\jdanswer{
$\frac{R^{2}g}{G}$
}
%% memo:
\memoanswer{
本题原卷及材料中未提供详解，待补充。
}

\item
%% number: 16
%% paperName: 1996年上海高考第 16 题 4 分
%% typeId: 3
%% type: 填空题
%% chapter: 磁场
%% point: 磁力矩
%% method: 无
%% score: 4
%% degree: 450
%% duplicateId: 0
%% body:
（B 组）如图所示，在光滑水平桌面上，有两根弯成直角的相同金属棒，它们的一端均可绕固定转动轴 O 自由转动，另一端 b 互相接触，组成一个正方形线框，每边长均为 $l$，匀强磁场的方向垂直桌面向下，磁感应强度为 $B$。当线框中通以图示方向的电流时，两金属棒在 b 点的相互作用力为 $f$，则此时线框中的电流大小为\tkanswer[2.4cm]{$\frac{\sqrt{2}f}{Bl}$}。（不计电流产生的磁场）
\onepicture[h=5.5cm, align=right]{\includesvg{figs/16}}
%% answer: 
\jdanswer{
$\frac{\sqrt{2}f}{Bl}$
}
%% memo:
\memoanswer{
本题原卷及材料中未提供详解，待补充。
}

\item
%% number: 17
%% paperName: 1996年上海高考第 17 题 4 分
%% typeId: 3
%% type: 填空题
%% chapter: 原子和原子核
%% point: 天然放射性
%% method: 无
%% score: 4
%% degree: 850
%% duplicateId: 0
%% body:
放射性元素 \ce{^{232}_{90}Th} 经过\tkanswer[1.2cm]{6}次 $\alpha$ 衰变和\tkanswer[1.2cm]{4}次 $\beta$ 衰变成为稳定元素 \ce{^{208}_{82}Pb}。
%% answer: 
\jdanswer{
6，4
}
%% memo:
\memoanswer{
本题原卷及材料中未提供详解，待补充。
}

\item
%% number: 18
%% paperName: 1996年上海高考第 18 题 4 分
%% typeId: 3
%% type: 填空题
%% chapter: 光学
%% point: 光的折射
%% method: 无
%% score: 4
%% degree: 650
%% duplicateId: 0
%% body:
如右图所示，有一个长方形容器，高为 $30\Ucm$，宽为 $40\Ucm$，在容器的底部平放着一把长 $40\Ucm$ 的刻度尺。眼睛在 OA 延长线上的 E 点观察，视线沿着 EA 斜向下看恰能看到尺的左端零刻度。现保持眼睛的位置不变，向容器内倒入某种液体且满至容器口，这时眼睛仍沿 EA 方向观察，恰能看到尺上 $20\Ucm$ 的刻度，则此种液体的折射率为\tkanswer[1.6cm]{1.44}。
\onepicture[w=5cm, align=right]{\includesvg{figs/18}}
%% answer: 
\jdanswer{
1.44（写成 $\frac{2\sqrt{13}}{5}$ 不扣分）
}
%% memo:
\memoanswer{
本题原卷及材料中未提供详解，待补充。
}

\item
%% number: 19
%% paperName: 1996年上海高考第 19 题 4 分
%% typeId: 3
%% type: 填空题
%% chapter: 力物体的平衡
%% point: 三力平衡
%% method: 无
%% score: 4
%% degree: 450
%% duplicateId: 0
%% body:
如右图所示，长为 $5\Um$ 的细绳的两端分别系于竖立在地面上相距为 $4\Um$ 的两杆的顶端 A、B。绳上挂一个光滑的轻质挂钩，其下连着一个重为 $12\UN$ 的物体。平衡时，绳中的张力 $T =$ \tkanswer[1.6cm]{10}$\UN$。
\onepicture[w=5cm, align=right]{\includesvg{figs/19}}
%% answer: 
\jdanswer{
$10\UN$
}
%% memo:
\memoanswer{
本题原卷及材料中未提供详解，待补充。
}

\item
%% number: 20
%% paperName: 1996年上海高考第 20 题 4 分
%% typeId: 3
%% type: 填空题
%% chapter: 交流电
%% point: LC振荡回路
%% method: 无
%% score: 4
%% degree: 650
%% duplicateId: 0
%% body:
如右图 LC 振荡回路中振荡电流的周期为 $2 \times 10^{-2}\Us$。自振荡电流沿反时针方向达最大值时开始计时，当 $t = 3.4 \times 10^{-2}\Us$ 时，电容器正处于\tkanswer[1.4cm]{充电}状态（填“充电”、“放电”、“充电完毕”或“放电完毕”）。这时电容器的上极板\tkanswer[1.6cm]{带正电}（填“带正电”、“带负电”或“不带电”）。
\onepicture[w=4cm, align=right]{\includesvg{figs/20}}
%% answer: 
\jdanswer{
充电，带正电
}
%% memo:
\memoanswer{
本题原卷及材料中未提供详解，待补充。
}

\item
%% number: 21
%% paperName: 1996年上海高考第 21 题 4 分
%% typeId: 3
%% type: 填空题
%% chapter: 运动和力
%% point: 动量守恒定律
%% method: 无
%% score: 4
%% degree: 450
%% duplicateId: 0
%% body:
总质量为 $M$ 的热气球由于故障在高空以匀速 $v$ 竖直下降。为了阻止继续下降，在 $t = 0$ 时刻，从热气球中释放了一个质量为 $m$ 的沙袋。不计空气阻力，当 $t =$ \tkanswer[2.4cm]{$\frac{(M-m)v}{mg}$} 时热气球停止下降，这时沙袋的速度为\tkanswer[1.8cm]{$\frac{M}{m}v$}（此时沙袋尚未着地）。
%% answer: 
\jdanswer{
$\frac{(M-m)v}{mg}$，$\frac{M}{m}v$
}
%% memo:
\memoanswer{
本题原卷及材料中未提供详解，待补充。
}

\item
%% number: 22
%% paperName: 1996年上海高考第 22 题 4 分
%% typeId: 4
%% type: 实验题
%% chapter: 光学
%% point: 光的干涉
%% method: 无
%% score: 4
%% degree: 650
%% duplicateId: 0
%% body:
用单色光做双缝干涉实验，下列说法中正确的是\xzanswer{A}
\fourchoices[answer=A]
{相邻干涉条纹之间的距离相等}
{中央明条纹宽度是两边明条纹宽度的 2 倍}
{屏与双缝之间距离减小，则屏上条纹间的距离增大}
{在实验装置不变的情况下，红光的条纹间距小于蓝光的条纹间距}
%% answer: 
\jdanswer{
A
}
%% memo:
\memoanswer{
本题原卷及材料中未提供详解，待补充。
}

\item
%% number: 23
%% paperName: 1996年上海高考第 23 题 7 分
%% typeId: 4
%% type: 实验题
%% chapter: 机械振动
%% point: 单摆测重力加速度
%% method: 无
%% score: 7
%% degree: 450
%% duplicateId: 0
%% body:
某同学在做“利用单摆测重力加速度”实验中，先测得摆线长为 $97.50\Ucm$，摆球直径为 $2.0\Ucm$，然后用秒表记录了单摆振动 50 次所用的时间（如图），则：
\begin{enumerate}
	\item
	该摆摆长为\tkanswer[1.8cm]{98.5}$\Ucm$，秒表所示读数为\tkanswer[1.8cm]{99.8}$\Us$。
	\item
	（单选题）如果他测得的 $g$ 值偏小，可能的原因是\xzanswer{B}
	\fourchoices[answer=B]
	{测摆线长时摆线拉得过紧}
	{摆线上端未牢固地系于悬点，振动中出现松动，使摆线长度增加了}
	{开始计时时，秒表过迟按下}
	{实验中误将 49 次全振动数为 50 次}
	\item
	为了提高实验精度，在实验中可改变几次摆长 $l$ 并测出相应的周期 $T$，从而得出一组对应的 $l$ 与 $T$ 的数据，再以 $l$ 为横坐标，$T^{2}$ 为纵坐标将所得数据连成直线（如右图），并求得该直线的斜率为 $k$，则重力加速度 $g =$ \tkanswer[2.2cm]{$\frac{4\pi^{2}}{k}$}（用 $k$ 表示）。
\end{enumerate}
\twopicture[labelA=1996上海23a, labelB=1996上海23b, w=6cm, align=right]
{\includesvg{figs/23a}}
{\includesvg{figs/23b}}
%% answer: 
\jdanswer{
\begin{enumerate}
	\item
	98.5，99.8

	\item
	B

	\item
	$\frac{4\pi^{2}}{k}$

\end{enumerate}
}
%% memo:
\memoanswer{
本题原卷及材料中未提供详解，待补充。
}

\item
%% number: 24
%% paperName: 1996年上海高考第 24 题 4 分
%% typeId: 4
%% type: 实验题
%% chapter: 电磁感应
%% point: 验证电磁感应实验
%% method: 无
%% score: 4
%% degree: 650
%% duplicateId: 0
%% body:
（多选题）用如图所示的装置研究电磁感应现象，在图示情况，当电键闭合瞬时，观察到电流表指针向右偏转，电键闭合一段时间后，为使电流表指针向左偏转，可采用的方法有\xzanswer{CD}
\onepicture[w=7.5cm, align=right]{\includesvg{figs/24}}
\fourchoices[answer=CD]
{将变阻器滑动头向右端滑动}
{将一软铁棒插入线圈 A 中}
{将线圈 A 从线圈 B 中提出}
{迅速断开电键}
%% answer: 
\jdanswer{
CD
}
%% memo:
\memoanswer{
本题原卷及材料中未提供详解，待补充。
}

\item
%% number: 25
%% paperName: 1996年上海高考第 25 题 5 分
%% typeId: 4
%% type: 实验题
%% chapter: 电路
%% point: 伏安法测电阻
%% method: 无
%% score: 5
%% degree: 650
%% duplicateId: 0
%% body:
某同学在做测定小灯泡功率的实验中得到如下一组 $U$ 和 $I$ 的数据：

\begin{center}
\begin{tabular}{|c|c|c|c|c|c|c|c|c|}
\hline
编号 & 1 & 2 & 3 & 4 & 5 & 6 & 7 & 8 \\ \hline
$U/\UV$ & 0.20 & 0.60 & 1.00 & 1.40 & 1.80 & 2.20 & 2.60 & 3.00 \\ \hline
$I/\UA$ & 0.020 & 0.060 & 0.100 & 0.140 & 0.170 & 0.190 & 0.200 & 0.205 \\ \hline
灯泡发光情况 & \multicolumn{2}{c|}{不亮} & \multicolumn{2}{c|}{微亮} & \multicolumn{2}{c|}{逐渐变亮} & \multicolumn{2}{c|}{正常发光} \\ \hline
\end{tabular}
\end{center}

\begin{enumerate}
	\item
	在右图上画出 $I-U$ 图线。
	\item
	从图线上可以看出，当功率逐渐增大时，灯丝电阻的变化情况是\tkanswer[3.5cm]{开始不变，后来逐渐增大}。
	\item
	这表明导体的电阻随温度升高而\tkanswer[1.4cm]{增大}。
\end{enumerate}
\onepicture[w=8cm, align=right]{\includesvg{figs/25}}
%% answer: 
\jdanswer{
\begin{enumerate}
	\item
	见图：
	\onepicture[w=8cm]{\includesvg{figs/25_an}}

	\item
	开始不变，后来逐渐增大

	\item
	增大

\end{enumerate}
}
%% memo:
\memoanswer{
本题原卷及材料中未提供详解，待补充。
}

\item
%% number: 26
%% paperName: 1996年上海高考第 26 题 6 分
%% typeId: 4
%% type: 实验题
%% chapter: 电路
%% point: 测量电阻率
%% method: 无
%% score: 6
%% degree: 450
%% duplicateId: 0
%% body:
如下图所示，P 是一根表面均匀地镀有很薄的发热电阻膜的长陶瓷管（其长度 $l$ 为 $50\Ucm$ 左右，直径 $D$ 为 $10\Ucm$ 左右），镀膜材料的电阻率 $\rho$ 已知，管的两端有导电箍 MN。现给你米尺，电压表 V、电流表 A、电源 $E$，滑动变阻器 $R$，电键 K 和若干导线，请设计一个测定膜层厚度 $d$ 的实验方案。
\begin{enumerate}
	\item
	实验中应测定的物理量是：\tkanswer[2.6cm]{$l$，$D$，$U$，$I$}。
	\item
	在右框内用符号画出测量电路图。
	\item
	计算膜层厚度的公式是：\tkanswer[2.8cm]{$d = \frac{l\rho I}{\pi DU}$}。
\end{enumerate}
\onepicture[w=5.5cm, align=right]{\includesvg{figs/26}}
%% answer: 
\jdanswer{
\begin{enumerate}
	\item
	$l$，$D$，$U$，$I$

	\item
	见图（电流表内接、外接均可）：
	\onepicture[w=6cm]{\includesvg{figs/26_an}}

	\item
	$d = \frac{l\rho I}{\pi DU}$

\end{enumerate}
}
%% memo:
\memoanswer{
本题原卷及材料中未提供详解，待补充。
}

\item
%% number: 27
%% paperName: 1996年上海高考第 27 题 10 分
%% typeId: 6
%% type: 计算题
%% chapter: 气体性质
%% point: 气体连接体
%% method: 无
%% score: 10
%% degree: 450
%% duplicateId: 0
%% body:
图示粗细均匀的 U 形管，右臂上端封闭，左臂中有一活塞，开始时用手握住活塞，使它与封闭端位于同一高度，这时两臂液面位于同一水平面内。管内液体的密度为 $\rho$，液体上方各有一定质量的理想气体，气柱长均为 $h$。今将活塞由图示的位置向上移动，移动的距离为 $2h$，这时两臂液面的高度差为 $h$。设整个过程中气体温度不变。问：活塞移动前，左右两臂液面上方气体的压强各为多少？
\onepicture[h=5.5cm, align=right]{\includesvg{figs/27}}
%% answer: 
\jdanswer{
$\frac{15}{4}\rho gh$
}
%% memo:
\memoanswer{
\onepicture[h=5.5cm, align=right]{\includesvg{figs/27_an}}

设截面积为 $S$，活塞移动前后左、右臂气体压强分别为 $p_{0\text{左}}$、$p_{\text{左}}$、$p_{0\text{右}}$、$p_{\text{右}}$，由玻意耳定律：

$p_{0\text{左}}hS = p_{\text{左}}\cdot 2.5hS$　①

$p_{0\text{右}}hS = p_{\text{右}}\cdot 1.5hS$　②

但 $p_{0\text{左}} = p_{0\text{右}}$　③

$p_{\text{右}} - p_{\text{左}} = \rho gh$　④

由①、②、③得 $2.5p_{\text{左}}h = 1.5p_{\text{右}}h$，

④代入上式得 $2.5p_{\text{左}} = 1.5(p_{\text{左}} + \rho gh)$，

得 $p_{\text{左}} = 1.5\rho gh$。

代入①式得 $p_{0\text{左}} = p_{0\text{右}} = 2.5p_{\text{左}} = \frac{5}{2}\cdot\frac{3}{2}\rho gh = \frac{15}{4}\rho gh$。

\textbf{评分标准}：列出①、②、④各 2 分，列出③ 1 分，得出结果 3 分。
}

\item
%% number: 28
%% paperName: 1996年上海高考第 28 题 12 分
%% typeId: 6
%% type: 计算题
%% chapter: 功和能
%% point: 整体机械能守恒
%% method: 无
%% score: 12
%% degree: 450
%% duplicateId: 0
%% body:
如图所示，半径为 $r$，质量不计的圆盘盘面与地面相垂直，圆心处有一个垂直盘面的光滑水平固定轴 O，在盘的最右边缘固定一个质量为 $m$ 的小球 A，在 O 点的正下方离 O 点 $\frac{r}{2}$ 处固定一个质量也为 $m$ 的小球 B。放开盘让其自由转动，问：
\begin{enumerate}
	\item
	当 A 球转到最低点时，两小球的重力势能之和减少了多少？
	\item
	A 球转到最低点时的线速度是多少？
	\item
	在转动过程中半径 OA 向左偏离竖直方向的最大角度是多少？
\end{enumerate}
\onepicture[w=4.5cm, align=right]{\includesvg{figs/28}}
%% answer: 
\jdanswer{
\begin{enumerate}
	\item
	$\frac{1}{2}mgr$

	\item
	$\sqrt{\frac{4gr}{5}}$

	\item
	$37^\circ$

\end{enumerate}
}
%% memo:
\memoanswer{
（1）$E_{p0} - E_{p} = mgr - mg\cdot\frac{r}{2} = \frac{1}{2}mgr$　①

（2）$\frac{1}{2}mgr = \frac{1}{2}mv_{A}^{2} + \frac{1}{2}m\left(\frac{v_{A}}{2}\right)^{2} = \frac{1}{2}m\cdot\frac{5}{4}v_{A}^{2}$　②

$v_{A} = \sqrt{\frac{4gr}{5}}$

\onepicture[w=4cm, align=right]{\includesvg{figs/28_an}}

（3）$mgr\cos\theta - mg\cdot\frac{r}{2}(1 + \sin\theta) = 0$　③

$2\cos\theta = 1 + \sin\theta$，

$4(1 - \sin^{2}\theta) = 1 + \sin^{2}\theta + 2\sin\theta$

$5\sin^{2}\theta + 2\sin\theta - 3 = 0$

$\sin\theta = \frac{-1 \pm \sqrt{1 + 15}}{5} = \frac{3}{5}$（舍去负根）

$\theta = 37^\circ$

\textbf{评分标准}：第（1）小题 3 分。第（2）小题 5 分，列出② 3 分，得出结果 2 分。第（3）小题 4 分，列出③ 2 分，得出 $\theta$ 结果 2 分。
}

\item
%% number: 29
%% paperName: 1996年上海高考第 29 题 13 分
%% typeId: 6
%% type: 计算题
%% chapter: 电路
%% point: 电容
%% method: 无
%% score: 13
%% degree: 450
%% duplicateId: 0
%% body:
三块相同的金属平板 A、B、D 自上而下水平放置，间距分别为 $h$ 和 $d$，如下图所示。A、B 两板中心开孔，在 A 板的开孔上搁有一金属容器 P，与 A 板接触良好，其内盛有导电液体。A 板通过闭合的电键 K 与电动势为 $U_{0}$ 的电池的正极相连，B 板与电池的负极相连并接地。容器 P 内的液体在底部小孔 O 处形成质量为 $m$、带电量为 $q$ 的液滴后自由下落，穿过 B 板的开孔 O$'$ 落在 D 板上，其电荷被 D 板吸附，液体随即蒸发，接着容器底部又形成相同的液滴自由下落，如此继续。设整个装置放在真空中。
\begin{enumerate}
	\item
	第 1 个液滴到达 D 板时的速度为多少？
	\item
	D 板最终可达到多高的电势？
	\item
	设液滴的电量是 A 板所带电量的 $\alpha$ 倍（$\alpha = 0.02$），A 板与 B 板构成的电容器的电容为 $C_{0} = 5 \times 10^{-12}\UF$，$U_{0} = 1000\UV$，$m = 0.02\Ug$，$h = d = 5\Ucm$。试计算 D 板最终的电势值；
	\item
	如果电键 K 不是始终闭合，而只是在第一个液滴形成前闭合一下，随即打开，其他条件与（3）相同。在这种情况下，D 板最终可达到的电势值为多少？说明理由。
\end{enumerate}
\onepicture[w=6.5cm, align=right]{\includesvg{figs/29}}
%% answer: 
\jdanswer{
\begin{enumerate}
	\item
	$\sqrt{2g(h + d) + \frac{2qU_{0}}{m}}$

	\item
	$U_{0} + \frac{mg(h + d)}{q}$

	\item
	$2.01 \times 10^{5}\UV$

	\item
	$1000\UV$

\end{enumerate}
}
%% memo:
\memoanswer{
（1）设第 1 个液滴到达 D 板时的速度为 $v$，由动能定理

$\frac{1}{2}mv^{2} = mg(h + d) + qU_{0}$　①

得 $v = \sqrt{2g(h + d) + \frac{2qU_{0}}{m}}$。

（2）当 D 板电势为 $U$ 时，液滴到达 D 板时的动能为

$E_{k} = mg(h + d) + qU_{0} - qU$　②

令 $E_{k} = 0$，得 $U = U_{0} + \frac{mg(h + d)}{q}$。

（3）$U = U_{0} + \frac{mg(h + d)}{\alpha C_{0}U_{0}} = 1000 + \frac{0.02 \times 10^{-3} \times 10 \times (0.05 + 0.05)}{0.02 \times 5 \times 10^{-12} \times 1000}\UV = 2.01 \times 10^{5}\UV$。

（4）$U$ 至多等于 A 板电荷全都到达 D 板时 D 板的电势值，由于 $h = d$，B、D 板间电容也是 $C_{0}$，故 $U$ 至多为 $U_{0}$，问题是 $U$ 能否达到 $U_{0}$。

由②得 $E_{k} = mg(h + d) + qU_{0} - qU > mg(h + d) - qU > mg(h + d) - q_{m}U_{0} = mg(h + d) - \alpha C_{0}U_{0}^{2}$，

其中 $q_{m} = \alpha C_{0}U_{0}$ 是 $q$ 的最大值，即第 1 个液滴的带电量。以（3）的数据代入，

$mg(h + d) - \alpha C_{0}U_{0}^{2} = 2 \times 10^{-5} - 0.02 \times 5 \times 10^{-12} \times 10^{6} > 0$。

可见恒有 $E_{k} > 0$，液滴一直往下滴，直至 A 板上电量全部转移到 D 板，$U = U_{0} = 1000\UV$。

\textbf{评分标准}：第（1）小题 2 分。列出①式 1 分，得出 $v$ 结果 1 分。第（2）小题 4 分。列出②式 2 分，由 $E_{k} = 0$ 得出 $U$ 结果 2 分。第（3）小题 3 分。知道 $q = \alpha C_{0}U_{0}$ 得 1 分，得出 $U$ 值 2 分。第（4）小题 4 分。正确论证 $E_{k}$ 恒大于零得 3 分，得出 $U = 1000\UV$ 得 1 分。
}

\end{enumerate}

\end{document}
